\subsection*{a) Modelo Oculto de Márkov}
A partir de los datos de entrenamiento:\\

Sea $\Sigma$, el conjunto de \textbf{observaciones}:
\[
\Sigma = \{0,\ 1\}
\]

Sea $S$, el conjunto de \textbf{estados ocultos} o \textbf{emisiones:}
\[
S = \{a,\ b\}
\]

Sea $\Pi$, la \textbf{matriz de probabilidades inciales}, la obtenemos al observar el \textbf{primer estado} de cada secuencia $y$ y contando cuántas veces inicia en cada estado:
\begin{center}
  \begin{tabular}{|c|c|}
    \hline
    Secuencia & Primer estado \\
    \hline
    aabb & $a$ \\
    bbaa & $b$ \\
    aabb & $a$ \\
    bbaa & $b$ \\
    \hline
  \end{tabular}
\end{center}

\[
\Pi_a = \frac{2}{4} = 0.5, \qquad \Pi_b = \frac{2}{4} = 0.5
\]
\[
\Pi = \begin{pmatrix} 0.5 \\ 0.5 \end{pmatrix}
\]

\bigskip

Sea $A$, la \textbf{matriz de probabilidades de transición}, la obtenemos al recorrer cada secuencia $y$ \textbf{par a par} (posición $t$ con posición $t+1$) y contando cuántas veces aparece cada transición consecutiva $s_i \to s_j$:

\begin{center}
  \begin{tabular}{|c|l|}
    \hline
    Secuencia $y$ & Pares consecutivos \\
    \hline
    \textbf{a\,a\,b\,b} & $(a{\to}a),\ (a{\to}b),\ (b{\to}b)$ \\
    \textbf{b\,b\,a\,a} & $(b{\to}b),\ (b{\to}a),\ (a{\to}a)$ \\
    \textbf{a\,a\,b\,b} & $(a{\to}a),\ (a{\to}b),\ (b{\to}b)$ \\
    \textbf{b\,b\,a\,a} & $(b{\to}b),\ (b{\to}a),\ (a{\to}a)$ \\
    \hline
  \end{tabular}
\end{center}

\noindent
Sumando los conteos de cada par sobre las cuatro secuencias:

\begin{center}
  \begin{tabular}{|c|c|c|c|c|c|}
    \hline
    Par & Seq.\,1 & Seq.\,2 & Seq.\,3 & Seq.\,4 & \textbf{Total} \\
    \hline
    $a \to a$ & 1 & 1 & 1 & 1 & \textbf{4} \\
    $a \to b$ & 1 & 0 & 1 & 0 & \textbf{2} \\
    $b \to a$ & 0 & 1 & 0 & 1 & \textbf{2} \\
    $b \to b$ & 1 & 1 & 1 & 1 & \textbf{4} \\
    \hline
  \end{tabular}
\end{center}

\noindent
Cada fila se normaliza diviendo entre el total de transiciones que parten
de ese estado ($4+2=6$ para $a$; $2+4=6$ para $b$):

\begin{center}
  \begin{tabular}{|c|c|c|}
    \hline
    & $a$ & $b$ \\
    \hline
    $a$ & 4 & 2\\
    $b$ & 2 & 4 \\
    \hline
  \end{tabular}
\end{center}

\[
A_{aa} = \frac{4}{6} = \frac{2}{3}, \quad
A_{ab} = \frac{2}{6} = \frac{1}{3}, \quad
A_{ba} = \frac{2}{6} = \frac{1}{3}, \quad
A_{bb} = \frac{4}{6} = \frac{2}{3}
\]
\[
A = \begin{pmatrix} 2/3 & 1/3 \\ 1/3 & 2/3 \end{pmatrix}
\approx
\begin{pmatrix} 0.667 & 0.333 \\ 0.333 & 0.667 \end{pmatrix}
\]

\bigskip

Sea $B$, la \textbf{matriz de probabilidades de emisión}, la obtenemos alineando \textbf{posición a posición} cada símbolo de $x_t$ con su estado oculto $y_t$ y contando cuántas veces cada estado emitió cada símbolo:
\begin{center}
  \begin{tabular}{|l|c|c|c|}
    \hline
    Secuencia & posición & $x_t$ & $y_t$ \\
    \hline
    $x{=}0101,\ y{=}aabb$ & 1 & 0 & $a$ \\
    & 2 & 1 & $a$ \\
    & 3 & 0 & $b$ \\
    & 4 & 1 & $b$ \\
    \hline
    $x{=}1010,\ y{=}bbaa$ & 1 & 1 & $b$ \\
    & 2 & 0 & $b$ \\
    & 3 & 1 & $a$ \\
    & 4 & 0 & $a$ \\
    \hline
    $x{=}0011,\ y{=}aabb$ & 1 & 0 & $a$ \\
    & 2 & 0 & $a$ \\
    & 3 & 1 & $b$ \\
    & 4 & 1 & $b$ \\
    \hline
    $x{=}1100,\ y{=}bbaa$ & 1 & 1 & $b$ \\
    & 2 & 1 & $b$ \\
    & 3 & 0 & $a$ \\
    & 4 & 0 & $a$ \\
    \hline
  \end{tabular}
\end{center}

Sumando:
\begin{center}
  \begin{tabular}{|c|c|c|c|}
    \hline
    Estado & emite $0$ & emite $1$ & Total \\
    \hline
    $a$ & 6 & 2 & 8 \\
    $b$ & 2 & 6 & 8 \\
    \hline
  \end{tabular}
\end{center}

Normalizando por fila (cada estado tiene 8 emisiones en total):
\begin{center}
  \begin{tabular}{|c|c|c|}
    \hline
    & $0$ & $1$ \\
    \hline
    $a$ & $6/8 = 0.75$ & $2/8 = 0.25$ \\
    $b$ & $2/8 = 0.25$ & $6/8 = 0.75$ \\
    \hline
  \end{tabular}
  \qquad
  $B = \begin{pmatrix} 0.75 & 0.25 \\ 0.25 & 0.75 \end{pmatrix}$
\end{center}

\bigskip

Por lo tanto el Modelo Oculto de Márkov es:
\[
HMM = (\Sigma,\ S,\ A,\ \Pi,\ B)
\]

\[
\Sigma = \{0, 1\}, \quad S = \{a, b\}
\]

\[
A = \begin{pmatrix} 2/3 & 1/3 \\ 1/3 & 2/3 \end{pmatrix},
\quad
\Pi = \begin{pmatrix} 0.5 \\ 0.5 \end{pmatrix},
\quad
B = \begin{pmatrix} 0.75 & 0.25 \\ 0.25 & 0.75 \end{pmatrix}
\]


\subsection*{b) Algoritmo de Viterbi}

La fórmula de \textbf{inicialización} usada será:
\[
\delta_i(1) = \pi_i \cdot B_{i,o_1}, \quad \forall\, i \in S
\]
La fórmula de \textbf{recursión} usada será:
\[
\delta_j(t) = \max_{i \in S}\bigl(\delta_i(t-1)\cdot A_{ij}\bigr)\cdot B_{j,o_t},
\qquad
\phi_j(t) = \arg\max_{i \in S}\bigl(\delta_i(t-1)\cdot A_{ij}\bigr)
\]

El \textbf{retroceso} identifica la cadena óptima:
\[
q^* = \arg\max_{i \in S}\,\delta_i(T), \qquad
\hat{y}^{(t)} = \phi_{\hat{y}^{(t+1)}}(t+1)
\]

\subsubsection*{1) Cadena de observación \texttt{00}:}

Diagrama de Trellis correspondiente:
\begin{center}
  \begin{tikzpicture}
    \def\xS{0} \def\xA{3.5} \def\xB{7}
    \def\ya{1.4} \def\yb{-1.4}
    
    \node[inicio] (ini) at (\xS,0) {Start};
    \node[font=\large] at (\xA, 3.0) {\texttt{o}$_1$ = 0};
    \node[font=\large] at (\xB, 3.0) {\texttt{o}$_2$ = 0};
    
    \node[state] (a1) at (\xA,\ya) {$a$};
    \node[state] (b1) at (\xA,\yb) {$b$};
    \node[state] (a2) at (\xB,\ya) {$a$};
    \node[state] (b2) at (\xB,\yb) {$b$};
    
    \draw[trans] (ini) -- (a1);
    \draw[trans] (ini) -- (b1);
    \draw[trans] (a1) -- (a2);
    \draw[trans] (a1) -- (b2);
    \draw[trans] (b1) -- (a2);
    \draw[trans] (b1) -- (b2);
  \end{tikzpicture}
\end{center}

\begin{enumerate}
  
\item \textbf{Inicialización} ($t=1$, observación: \texttt{0})
  
  No hay estado anterior; se usa la distribución inicial:
  \[
  \delta_a(1) = \Pi_a \cdot B_{a,0} = 0.5 \times 0.75 = \mathbf{0.375}
  \]
  \[
  \delta_b(1) = \Pi_b \cdot B_{b,0} = 0.5 \times 0.25 = \mathbf{0.125}
  \]
  No se calcula $\phi$ en $t=1$.
  
  \begin{center}
    \begin{tikzpicture}
      \def\xS{0} \def\xA{4} \def\xB{7.5}
      \def\ya{1.4} \def\yb{-1.4}
      
      \node[inicio] (ini) at (\xS,0) {Start};
      \node[font=\large] at (\xA, 2.8) {\texttt{o}$_1$ = 0};
      \node[font=\large] at (\xB, 2.8) {\texttt{o}$_2$ = 0};
      
      \node[state] (a1) at (\xA,\ya) {$a$};
      \node[state] (b1) at (\xA,\yb) {$b$};
      \node[state] (a2) at (\xB,\ya) {$a$};
      \node[state] (b2) at (\xB,\yb) {$b$};
      
      \draw[winner] (ini) -- node[above, yshift=4pt]  {$0.375$} (a1);
      \draw[winner] (ini) -- node[below, yshift=-4pt] {$0.125$} (b1);
      \draw[trans] (a1)--(a2); \draw[trans] (a1)--(b2);
      \draw[trans] (b1)--(a2); \draw[trans] (b1)--(b2);
    \end{tikzpicture}
  \end{center}
  
\item \textbf{Recursión} ($t=2$, observación: \texttt{0})
  
  $\delta_a(2)$, llegar a $a$:
  \[
  \delta_a(1)\cdot A_{aa} = 0.375 \times \tfrac{2}{3} = 0.250
  \]
  \[
  \delta_b(1)\cdot A_{ba} = 0.125 \times \tfrac{1}{3} = 0.0417
  \]
  \[
  \max(0.250,\ 0.0417) = 0.250 \quad\Rightarrow\quad \phi_a(2) = a
  \]
  \[
  \delta_a(2) = 0.250 \times B_{a,0} = 0.250 \times 0.75 = \mathbf{0.1875}
  \]
  
  $\delta_b(2)$, llegar a $b$:
  \[
  \delta_a(1)\cdot A_{ab} = 0.375 \times \tfrac{1}{3} = 0.125
  \]
  \[
  \delta_b(1)\cdot A_{bb} = 0.125 \times \tfrac{2}{3} = 0.0833
  \]
  \[
  \max(0.125,\ 0.0833) = 0.125 \quad\Rightarrow\quad \phi_b(2) = a
  \]
  \[
  \delta_b(2) = 0.125 \times B_{b,0} = 0.125 \times 0.25 = \mathbf{0.03125}
  \]
  
  \begin{center}
    \begin{tikzpicture}
      \def\xS{-1} \def\xA{3.5} \def\xB{8.5}
      \def\ya{1.4} \def\yb{-1.4}
      
      \node[inicio] (ini) at (\xS,0) {Start};
      \node[font=\large] at (\xA, 2.8) {\texttt{o}$_1$ = 0};
      \node[font=\large] at (\xB, 2.8) {\texttt{o}$_2$ = 0};
      
      \node[ellipstate] (a1) at (\xA,\ya) {$\delta_a(1)=0.375$};
      \node[ellipstate] (b1) at (\xA,\yb) {$\delta_b(1)=0.125$};
      \node[state] (a2) at (\xB,\ya) {$a$};
      \node[state] (b2) at (\xB,\yb) {$b$};
      
      \draw[trans] (ini)--(a1); \draw[trans] (ini)--(b1);
      % ganadoras (azul): ambas desde a1
      \draw[winner] (a1) -- node[above, yshift=4pt]  {$0.1875$} (a2);
      \draw[winner] (a1) -- node[below, yshift=-4pt] {$0.03125$} (b2);
      \draw[loser]  (b1) -- (a2);
      \draw[loser]  (b1) -- (b2);
    \end{tikzpicture}
  \end{center}
  
\item \textbf{Retroceso (Backtracking)}
  
  \[
  \max\bigl(\delta_a(2),\ \delta_b(2)\bigr)
  = \max(0.1875,\ 0.03125) = 0.1875
  \;\Rightarrow\; \hat{y}^{(2)} = a
  \]
  \[
  t=2:\ a \xrightarrow{\phi_a(2)=a} t=1:\ a
  \]
  
  La cadena de emisiones más probable para \texttt{00} es:
  \[
  \boxed{\hat{y}^{(1)}\hat{y}^{(2)} = a \to a}
  \]
  con probabilidad máxima $\mathbf{0.1875}$.
  
  \begin{center}
    \begin{tikzpicture}
      \def\xS{-1} \def\xA{3.5} \def\xB{8.5}
      \def\ya{1.4} \def\yb{-1.4}
      
      \node[inicio, draw=red!70, fill=red!5] (ini) at (\xS,0) {Start};
      \node[font=\large] at (\xA, 2.8) {\texttt{o}$_1$ = 0};
      \node[font=\large] at (\xB, 2.8) {\texttt{o}$_2$ = 0};
      
      \node[ellipstate, draw=red!70, fill=red!5] (a1) at (\xA,\ya) {$\delta_a(1)=0.375$};
      \node[ellipstate] (b1) at (\xA,\yb) {$\delta_b(1)=0.125$};
      \node[ellipstate, draw=red!70, fill=red!5] (a2) at (\xB,\ya) {$\delta_a(2)=\mathbf{0.1875}$};
      \node[ellipstate] (b2) at (\xB,\yb) {$\delta_b(2)=0.03125$};
      
      \draw[loser]   (ini)--(b1);
      \draw[optimal] (ini)--(a1);
      \draw[optimal] (a1) --(a2);
      \draw[loser]   (a1) --(b2);
      \draw[loser]   (b1) --(a2);
      \draw[loser]   (b1) --(b2);
    \end{tikzpicture}
  \end{center}
  
\end{enumerate}

\subsubsection*{2) Cadena de observación \texttt{01}:}
Con esto, damos el siguiente diagrama de Trellis correspondiente:

\begin{center}
  \begin{tikzpicture}
    \def\xS{0} \def\xA{3.5} \def\xB{7}
    \def\ya{1.4} \def\yb{-1.4}
    
    \node[inicio] (ini) at (\xS,0) {Start};
    \node[font=\large] at (\xA, 3.0) {\texttt{o}$_1$ = 0};
    \node[font=\large] at (\xB, 3.0) {\texttt{o}$_2$ = 1};
    
    \node[state] (a1) at (\xA,\ya) {$a$};
    \node[state] (b1) at (\xA,\yb) {$b$};
    \node[state] (a2) at (\xB,\ya) {$a$};
    \node[state] (b2) at (\xB,\yb) {$b$};
    
    \draw[trans] (ini)--(a1); \draw[trans] (ini)--(b1);
    \draw[trans] (a1)--(a2); \draw[trans] (a1)--(b2);
    \draw[trans] (b1)--(a2); \draw[trans] (b1)--(b2);
  \end{tikzpicture}
\end{center}

\begin{enumerate}
  
\item \textbf{Inicialización} ($t=1$, observación: \texttt{0})
  
  Igual que en el caso anterior (misma primera observación):
  \[
  \delta_a(1) = 0.5 \times 0.75 = \mathbf{0.375}, \qquad
  \delta_b(1) = 0.5 \times 0.25 = \mathbf{0.125}
  \]
  
  \begin{center}
    \begin{tikzpicture}
      \def\xS{0} \def\xA{4} \def\xB{7.5}
      \def\ya{1.4} \def\yb{-1.4}
      
      \node[inicio] (ini) at (\xS,0) {Start};
      \node[font=\large] at (\xA, 2.8) {\texttt{o}$_1$ = 0};
      \node[font=\large] at (\xB, 2.8) {\texttt{o}$_2$ = 1};
      
      \node[state] (a1) at (\xA,\ya) {$a$};
      \node[state] (b1) at (\xA,\yb) {$b$};
      \node[state] (a2) at (\xB,\ya) {$a$};
      \node[state] (b2) at (\xB,\yb) {$b$};
      
      \draw[winner] (ini) -- node[above, yshift=4pt]  {$0.375$} (a1);
      \draw[winner] (ini) -- node[below, yshift=-4pt] {$0.125$} (b1);
      \draw[trans] (a1)--(a2); \draw[trans] (a1)--(b2);
      \draw[trans] (b1)--(a2); \draw[trans] (b1)--(b2);
    \end{tikzpicture}
  \end{center}
  
\item \textbf{Recursión} ($t=2$, observación: \texttt{1})
  
  $\delta_a(2)$, llegar a $a$:
  \[
  \delta_a(1)\cdot A_{aa} = 0.375 \times \tfrac{2}{3} = 0.250
  \]
  \[
  \delta_b(1)\cdot A_{ba} = 0.125 \times \tfrac{1}{3} = 0.0417
  \]
  \[
  \max(0.250,\ 0.0417) = 0.250 \quad\Rightarrow\quad \phi_a(2) = a
  \]
  \[
  \delta_a(2) = 0.250 \times B_{a,1} = 0.250 \times 0.25 = \mathbf{0.0625}
  \]
  
  $\delta_b(2)$, llegar a $b$:
  \[
  \delta_a(1)\cdot A_{ab} = 0.375 \times \tfrac{1}{3} = 0.125
  \]
  \[
  \delta_b(1)\cdot A_{bb} = 0.125 \times \tfrac{2}{3} = 0.0833
  \]
  \[
  \max(0.125,\ 0.0833) = 0.125 \quad\Rightarrow\quad \phi_b(2) = a
  \]
  \[
  \delta_b(2) = 0.125 \times B_{b,1} = 0.125 \times 0.75 = \mathbf{0.09375}
  \]
  
  \begin{center}
    \begin{tikzpicture}
      \def\xS{-1} \def\xA{3.5} \def\xB{8.5}
      \def\ya{1.4} \def\yb{-1.4}
      
      \node[inicio] (ini) at (\xS,0) {Start};
      \node[font=\large] at (\xA, 2.8) {\texttt{o}$_1$ = 0};
      \node[font=\large] at (\xB, 2.8) {\texttt{o}$_2$ = 1};
      
      \node[ellipstate] (a1) at (\xA,\ya) {$\delta_a(1)=0.375$};
      \node[ellipstate] (b1) at (\xA,\yb) {$\delta_b(1)=0.125$};
      \node[state] (a2) at (\xB,\ya) {$a$};
      \node[state] (b2) at (\xB,\yb) {$b$};
      
      \draw[trans] (ini)--(a1); \draw[trans] (ini)--(b1);
      \draw[winner] (a1) -- node[above, yshift=4pt]  {$0.0625$}  (a2);
      \draw[winner] (a1) -- node[below, yshift=-4pt] {$0.09375$} (b2);
      \draw[loser]  (b1) -- (a2);
      \draw[loser]  (b1) -- (b2);
    \end{tikzpicture}
  \end{center}
  
\item \textbf{Retroceso (Backtracking)}
  
  \[
  \max\bigl(\delta_a(2),\ \delta_b(2)\bigr)
  = \max(0.0625,\ 0.09375) = 0.09375
  \;\Rightarrow\; \hat{y}^{(2)} = b
  \]
  \[
  t=2:\ b \xrightarrow{\phi_b(2)=a} t=1:\ a
  \]
  
La cadena de emisiones más probable para \texttt{01} es:
\[
\boxed{\hat{y}^{(1)}\hat{y}^{(2)} = a \to b}
\]
con probabilidad máxima $\mathbf{0.09375}$.
 
\begin{center}
\begin{tikzpicture}
  \def\xS{-1} \def\xA{3.5} \def\xB{8.5}
  \def\ya{1.4} \def\yb{-1.4}
 
  \node[inicio, draw=red!70, fill=red!5] (ini) at (\xS,0) {Start};
  \node[font=\large] at (\xA, 2.8) {\texttt{o}$_1$ = 0};
  \node[font=\large] at (\xB, 2.8) {\texttt{o}$_2$ = 1};
 
  \node[ellipstate, draw=red!70, fill=red!5] (a1) at (\xA,\ya) {$\delta_a(1)=0.375$};
  \node[ellipstate] (b1) at (\xA,\yb) {$\delta_b(1)=0.125$};
  \node[ellipstate] (a2) at (\xB,\ya) {$\delta_a(2)=0.0625$};
  \node[ellipstate, draw=red!70, fill=red!5] (b2) at (\xB,\yb) {$\delta_b(2)=\mathbf{0.09375}$};
 
  \draw[loser]   (ini)--(b1);
  \draw[optimal] (ini)--(a1);
  \draw[loser]   (a1) --(a2);
  \draw[optimal] (a1) --(b2);
  \draw[loser]   (b1) --(a2);
  \draw[loser]   (b1) --(b2);
\end{tikzpicture}
\end{center}
 
\end{enumerate}
 
\subsubsection*{3) Cadena de observación \texttt{010}:}
Con esto, damos el siguiente diagrama de Trellis correspondiente:
 
\begin{center}
\begin{tikzpicture}
  \def\xS{0} \def\xA{3} \def\xB{6} \def\xC{9}
  \def\ya{1.4} \def\yb{-1.4}
 
  \node[inicio] (ini) at (\xS,0) {Start};
  \node[font=\normalsize] at (\xA, 2.8) {\texttt{o}$_1$=0};
  \node[font=\normalsize] at (\xB, 2.8) {\texttt{o}$_2$=1};
  \node[font=\normalsize] at (\xC, 2.8) {\texttt{o}$_3$=0};
 
  \node[state] (a1) at (\xA,\ya) {$a$};
  \node[state] (b1) at (\xA,\yb) {$b$};
  \node[state] (a2) at (\xB,\ya) {$a$};
  \node[state] (b2) at (\xB,\yb) {$b$};
  \node[state] (a3) at (\xC,\ya) {$a$};
  \node[state] (b3) at (\xC,\yb) {$b$};
 
  \draw[trans] (ini)--(a1); \draw[trans] (ini)--(b1);
  \draw[trans] (a1)--(a2); \draw[trans] (a1)--(b2);
  \draw[trans] (b1)--(a2); \draw[trans] (b1)--(b2);
  \draw[trans] (a2)--(a3); \draw[trans] (a2)--(b3);
  \draw[trans] (b2)--(a3); \draw[trans] (b2)--(b3);
\end{tikzpicture}
\end{center}
 
\begin{enumerate}
 
\item \textbf{Inicialización} ($t=1$, observación: \texttt{0})
 
\[
  \delta_a(1) = 0.5 \times 0.75 = \mathbf{0.375}, \qquad
  \delta_b(1) = 0.5 \times 0.25 = \mathbf{0.125}
\]
 
\begin{center}
\begin{tikzpicture}
  \def\xS{-1} \def\xA{3.5} \def\xB{7} \def\xC{10.5}
  \def\ya{1.4} \def\yb{-1.4}
 
  \node[inicio] (ini) at (\xS,0) {Start};
  \node[font=\normalsize] at (\xA, 2.8) {\texttt{o}$_1$=0};
  \node[font=\normalsize] at (\xB, 2.8) {\texttt{o}$_2$=1};
  \node[font=\normalsize] at (\xC, 2.8) {\texttt{o}$_3$=0};
 
  \node[state] (a1) at (\xA,\ya) {$a$};
  \node[state] (b1) at (\xA,\yb) {$b$};
  \node[state] (a2) at (\xB,\ya) {$a$};
  \node[state] (b2) at (\xB,\yb) {$b$};
  \node[state] (a3) at (\xC,\ya) {$a$};
  \node[state] (b3) at (\xC,\yb) {$b$};
 
  \draw[winner] (ini)--node[above,yshift=4pt]{$0.375$}(a1);
  \draw[winner] (ini)--node[below,yshift=-4pt]{$0.125$}(b1);
  \draw[trans](a1)--(a2);\draw[trans](a1)--(b2);
  \draw[trans](b1)--(a2);\draw[trans](b1)--(b2);
  \draw[trans](a2)--(a3);\draw[trans](a2)--(b3);
  \draw[trans](b2)--(a3);\draw[trans](b2)--(b3);
\end{tikzpicture}
\end{center}
 
\item \textbf{Recursión} ($t=2$, observación: \texttt{1})
 
$\delta_a(2)$, llegar a $a$:
\[
  \delta_a(1)\cdot A_{aa} = 0.375\times\tfrac{2}{3} = 0.250
  \qquad
  \delta_b(1)\cdot A_{ba} = 0.125\times\tfrac{1}{3} = 0.0417
\]
\[
  \max(0.250,\ 0.0417)=0.250 \;\Rightarrow\; \phi_a(2)=a
  \qquad
  \delta_a(2)=0.250\times0.25=\mathbf{0.0625}
\]
 
$\delta_b(2)$, llegar a $b$:
\[
  \delta_a(1)\cdot A_{ab} = 0.375\times\tfrac{1}{3} = 0.125
  \qquad
  \delta_b(1)\cdot A_{bb} = 0.125\times\tfrac{2}{3} = 0.0833
\]
\[
  \max(0.125,\ 0.0833)=0.125 \;\Rightarrow\; \phi_b(2)=a
  \qquad
  \delta_b(2)=0.125\times0.75=\mathbf{0.09375}
\]
 
\begin{center}
\begin{tikzpicture}
  \def\xS{-2} \def\xA{2.5} \def\xB{7} \def\xC{11}
  \def\ya{1.4} \def\yb{-1.4}
 
  \node[inicio] (ini) at (\xS,0) {Start};
  \node[font=\normalsize] at (\xA, 2.8) {\texttt{o}$_1$=0};
  \node[font=\normalsize] at (\xB, 2.8) {\texttt{o}$_2$=1};
  \node[font=\normalsize] at (\xC, 2.8) {\texttt{o}$_3$=0};
 
  \node[ellipstate] (a1) at (\xA,\ya) {$\delta_a(1)=0.375$};
  \node[ellipstate] (b1) at (\xA,\yb) {$\delta_b(1)=0.125$};
  \node[state] (a2) at (\xB,\ya) {$a$};
  \node[state] (b2) at (\xB,\yb) {$b$};
  \node[state] (a3) at (\xC,\ya) {$a$};
  \node[state] (b3) at (\xC,\yb) {$b$};
 
  \draw[trans](ini)--(a1);\draw[trans](ini)--(b1);
  \draw[winner](a1)--node[above,yshift=4pt]{$0.0625$}(a2);
  \draw[winner](a1)--node[below,yshift=-4pt]{$0.09375$}(b2);
  \draw[loser](b1)--(a2);\draw[loser](b1)--(b2);
  \draw[trans](a2)--(a3);\draw[trans](a2)--(b3);
  \draw[trans](b2)--(a3);\draw[trans](b2)--(b3);
\end{tikzpicture}
\end{center}
 
\item \textbf{Recursión} ($t=3$, observación: \texttt{0})
 
$\delta_a(3)$, llegar a $a$:
\[
  \delta_a(2)\cdot A_{aa} = 0.0625\times\tfrac{2}{3} = 0.04167
\]
\[
  \delta_b(2)\cdot A_{ba} = 0.09375\times\tfrac{1}{3} = 0.03125
\]
\[
  \max(0.04167,\ 0.03125)=0.04167 \;\Rightarrow\; \phi_a(3)=a
\]
\[
  \delta_a(3)=0.04167\times B_{a,0}=0.04167\times0.75=\mathbf{0.03125}
\]
 
$\delta_b(3)$, llegar a $b$:
\[
  \delta_a(2)\cdot A_{ab} = 0.0625\times\tfrac{1}{3} = 0.02083
\]
\[
  \delta_b(2)\cdot A_{bb} = 0.09375\times\tfrac{2}{3} = 0.06250
\]
\[
  \max(0.02083,\ 0.06250)=0.06250 \;\Rightarrow\; \phi_b(3)=b
\]
\[
  \delta_b(3)=0.06250\times B_{b,0}=0.06250\times0.25=\mathbf{0.015625}
\]
 
\begin{center}
\begin{tikzpicture}
  \def\xS{-2} \def\xA{2.5} \def\xB{7} \def\xC{11.5}
  \def\ya{1.4} \def\yb{-1.4}
 
  \node[inicio] (ini) at (\xS,0) {Start};
  \node[font=\normalsize] at (\xA, 2.8) {\texttt{o}$_1$=0};
  \node[font=\normalsize] at (\xB, 2.8) {\texttt{o}$_2$=1};
  \node[font=\normalsize] at (\xC, 2.8) {\texttt{o}$_3$=0};
 
  \node[ellipstate] (a1) at (\xA,\ya) {$\delta_a(1)=0.375$};
  \node[ellipstate] (b1) at (\xA,\yb) {$\delta_b(1)=0.125$};
  \node[ellipstate] (a2) at (\xB,\ya) {$\delta_a(2)=0.0625$};
  \node[ellipstate] (b2) at (\xB,\yb) {$\delta_b(2)=0.09375$};
  \node[state] (a3) at (\xC,\ya) {$a$};
  \node[state] (b3) at (\xC,\yb) {$b$};
 
  \draw[trans](ini)--(a1);\draw[trans](ini)--(b1);
  \draw[trans](a1)--(a2);\draw[trans](a1)--(b2);
  \draw[loser](b1)--(a2);\draw[loser](b1)--(b2);
  % ganadoras t=3: a gana desde a2, b gana desde b2
  \draw[loser] (a2)--(a3);
  \draw[loser] (a2)--(b3);
  \draw[winner] (b2)--node[above,yshift=8pt]{$0.03125$}(a3);
  \draw[winner](b2)--node[below,yshift=-4pt]{$0.015625$}(b3);
\end{tikzpicture}
\end{center}
 
\item \textbf{Retroceso (Backtracking)}
 
\[
  \max\bigl(\delta_a(3),\ \delta_b(3)\bigr)
  = \max(0.03125,\ 0.015625) = 0.03125
  \;\Rightarrow\; \hat{y}^{(3)} = a
\]
\[
  t=3:\ a \xrightarrow{\phi_a(3)=a} t=2:\ a \xrightarrow{\phi_a(2)=a} t=1:\ a
\]
 
La cadena de emisiones más probable para \texttt{010} es:
\[
  \boxed{\hat{y}^{(1)}\hat{y}^{(2)}\hat{y}^{(3)} = a \to a \to a}
\]
con probabilidad máxima $\mathbf{0.03125}$.
 
\begin{tikzpicture}
  \def\xS{-6} \def\xA{-3} \def\xB{2} \def\xC{7.2}
  \def\ya{1.4} \def\yb{-1.4}
 
  \node[inicio, draw=red!70, fill=red!5] (ini) at (\xS,0) {Start};
  \node[font=\normalsize] at (\xA, 2.8) {\texttt{o}$_1$=0};
  \node[font=\normalsize] at (\xB, 2.8) {\texttt{o}$_2$=1};
  \node[font=\normalsize] at (\xC, 2.8) {\texttt{o}$_3$=0};
 
  \node[ellipstate, draw=red!70, fill=red!5] (a1) at (\xA,\ya) {$\delta_a(1)=0.375$};
  \node[ellipstate] (b1) at (\xA,\yb) {$\delta_b(1)=0.125$};
  \node[ellipstate] (a2) at (\xB,\ya) {$\delta_a(2)=0.0625$};
  \node[ellipstate, draw=red!70, fill=red!5] (b2) at (\xB,\yb) {$\delta_b(2)=0.09375$};
  \node[ellipstate, draw=red!70, fill=red!5] (a3) at (\xC,\ya) {$\delta_a(3)=\mathbf{0.03125}$};
  \node[ellipstate] (b3) at (\xC,\yb) {$\delta_b(3)=0.015625$};
 
  \draw[loser]   (ini)--(b1);
  \draw[optimal] (ini)--(a1);
  \draw[loser]   (a1) --(a2);
  \draw[optimal] (a1) --(b2);
  \draw[loser]   (b1) --(a2);
  \draw[loser]   (b1) --(b2);
  \draw[loser]   (a2) --(a3);
  \draw[loser]   (a2) --(b3);
  \draw[optimal] (b2) --(a3);
  \draw[loser]   (b2) --(b3);
\end{tikzpicture}x
 
\end{enumerate}
